% !TeX root = ../main.tex
% Add the above to each chapter to make compiling the PDF easier in some editors.

\chapter{Results and Evaluation}\label{chapter:results}

This chapter presents and interprets the outcomes of the experiments described in \autoref{chapter:experiments}.
Each section addresses one research dimension: image mode choice, feature engineering, single- vs.\ dual-channel, image resolution, and ensembling,
before placing the best results in context with the MM-AutoSolver baseline.

\section{Overview of All Experimental Results}\label{sec:results-overview}

Each of the following sections presents its results table alongside the discussion, introducing metrics as they become relevant.
Bold marks the best value in each column; for Fail\% and MRT$\times$, lower is better.

\section{Impact of Image Mode}\label{sec:results-image-mode}

% ── Campaign 1: Single-channel, 14 features, no penalty ──────────────────────
\begin{table}[H]
  \centering
  \caption{Campaign~1 — Single-channel: 14 features, no convergence penalty, small model.}
  \label{tab:results-campaign1}
  \resizebox{\textwidth}{!}{%
  \begin{tabular}{lrrrrrrrr}
    \toprule
    \textbf{Experiment} & \textbf{Acc\%} & \textbf{MP\%} & \textbf{MR\%} & \textbf{F1\%} & \textbf{Top-2\%} & \textbf{Top-3\%} & \textbf{Fail\%} & \textbf{MRT$\times$} \\
    \midrule
    1-NN                     & 60.06 & 58.94 & 58.56 & 57.97 & 82.04 & 90.30 & 2.27 & 1.533 \\
    5-NN                     & 62.95 & 62.46 & 61.85 & 61.12 & 82.56 & 90.40 & 3.20 & 1.605 \\
    nocnn                    & 58.31 & 60.66 & 61.14 & 57.57 & 78.95 & 86.38 & 5.88 & 1.345 \\
    \midrule
    binary\_64               & 59.55 & 63.33 & 62.69 & 58.80 & 82.04 & 90.61 & 6.30 & 1.274 \\
    binary\_128              & 60.06 & 63.65 & 62.95 & 59.26 & 83.28 & 91.95 & 5.47 & 2.819 \\
    density\_64              & 60.27 & 63.17 & 62.29 & 59.46 & 81.42 & 91.54 & 5.47 & 1.291 \\
    density\_128             & 63.26 & 65.25 & 65.57 & 62.25 & 84.42 & 93.19 & 4.44 & 1.271 \\
    diagonal\_64             & 59.75 & 62.92 & 62.82 & 58.79 & 82.35 & 90.92 & 5.37 & 4.196 \\
    diagonal\_128            & 59.44 & 63.07 & 62.53 & 58.79 & 81.63 & 90.51 & 5.99 & 1.316 \\
    log\_density\_64         & 62.33 & 65.39 & 64.86 & 61.38 & 82.97 & 92.88 & 5.16 & 2.807 \\
    log\_density\_128        & 62.64 & 64.33 & 64.98 & 61.69 & 83.90 & 92.67 & 5.16 & 1.271 \\
    \textbf{magnitude\_64}   & \textbf{64.19} & 66.19 & \textbf{66.55} & \textbf{63.42} & \textbf{86.58} & 93.91 & 3.92 & 1.259 \\
    magnitude\_128           & 63.88 & \textbf{67.58} & 66.35 & 62.99 & 86.48 & 93.81 & 3.92 & 1.115 \\
    sign\_64                 & 60.37 & 63.08 & 62.59 & 59.29 & 82.25 & 91.64 & 5.68 & 1.288 \\
    sign\_128                & 61.20 & 65.11 & 64.28 & 60.53 & 83.08 & 92.16 & 4.02 & 4.139 \\
    signed\_magnitude\_64    & 62.85 & 64.56 & 65.20 & 61.68 & 84.83 & 94.22 & 4.75 & \textbf{1.069} \\
    signed\_magnitude\_128   & 61.82 & 64.75 & 64.53 & 60.42 & 84.73 & 93.09 & 5.06 & 1.222 \\
    symmetry\_64             & 62.95 & 64.40 & 65.66 & 61.85 & 85.66 & 94.01 & \textbf{3.51} & 4.081 \\
    symmetry\_128            & 62.95 & 64.14 & 65.32 & 61.62 & 85.96 & \textbf{94.43} & 4.64 & 1.233 \\
    \midrule
    rcm\_binary\_64          & 60.06 & 62.61 & 62.84 & 59.17 & 82.46 & 91.64 & 6.40 & 1.322 \\
    rcm\_binary\_128         & 59.86 & 62.18 & 62.40 & 59.01 & 82.35 & 91.64 & 6.81 & 1.269 \\
    rcm\_density\_64         & 61.40 & 64.53 & 64.09 & 60.43 & 83.59 & 93.40 & 5.68 & 1.235 \\
    rcm\_density\_128        & 62.33 & 65.47 & 64.84 & 61.33 & 82.77 & 93.29 & 4.75 & 1.262 \\
    rcm\_log\_density\_64    & 62.85 & 65.78 & 65.32 & 61.88 & 83.90 & 93.29 & 4.75 & 2.770 \\
    rcm\_log\_density\_128   & 61.82 & 63.55 & 64.19 & 60.94 & 83.49 & 92.36 & 5.47 & 1.271 \\
    rcm\_magnitude\_64       & 62.85 & 65.97 & 65.50 & 62.14 & 84.83 & 93.19 & 4.64 & 1.288 \\
    rcm\_magnitude\_128      & 63.05 & 65.93 & 65.35 & 62.55 & 84.62 & 92.67 & 4.13 & 1.281 \\
    rcm\_sign\_64            & 61.71 & 65.19 & 64.23 & 60.84 & 82.56 & 92.47 & 4.64 & 2.802 \\
    rcm\_sign\_128           & 61.82 & 65.53 & 64.63 & 61.05 & 82.77 & 92.98 & 4.23 & 2.829 \\
    rcm\_signed\_magnitude\_64  & 63.67 & 66.12 & 65.51 & 62.24 & 85.66 & 94.01 & 4.02 & 2.586 \\
    rcm\_signed\_magnitude\_128 & 63.36 & 64.74 & 65.65 & 62.23 & 85.35 & 94.12 & 4.64 & 1.235 \\
    \bottomrule
  \end{tabular}%
  }
\end{table}

% ── Campaign 2: Single-channel, 20 features, convergence penalty λ=1.0 ────────
\begin{table}[H]
  \centering
  \caption{Campaign~2 — Single-channel: 20 features, convergence penalty $\lambda=1.0$, large model, 64\,px.}
  \label{tab:results-campaign2}
  \resizebox{\textwidth}{!}{%
  \begin{tabular}{lrrrrrrrr}
    \toprule
    \textbf{Experiment} & \textbf{Acc\%} & \textbf{MP\%} & \textbf{MR\%} & \textbf{F1\%} & \textbf{Top-2\%} & \textbf{Top-3\%} & \textbf{Fail\%} & \textbf{MRT$\times$} \\
    \midrule
    1-NN                        & 59.44 & 59.02 & 58.34 & 57.00 & 82.15 & 90.40 & 1.86 & 1.539 \\
    5-NN                        & 63.57 & 63.25 & 62.38 & 61.61 & 82.97 & 91.02 & 2.89 & 1.589 \\
    nocnn                       & 59.75 & 62.40 & 62.72 & 59.08 & 80.80 & 89.78 & 4.95 & 2.926 \\
    \midrule
    binary\_64                  & 59.34 & 62.66 & 62.53 & 58.87 & 81.53 & 91.33 & 5.06 & 1.343 \\
    density\_64                 & 52.63 & 55.76 & 54.56 & 50.74 & 72.14 & 79.26 & 5.16 & 1.363 \\
    diagonal\_64                & 59.13 & 62.96 & 62.11 & 58.46 & 81.63 & 90.61 & 5.16 & 1.371 \\
    log\_density\_64            & 62.23 & 66.12 & 65.13 & 61.65 & 82.15 & 92.36 & 4.54 & 1.284 \\
    \textbf{magnitude\_64}      & \textbf{64.29} & 66.58 & \textbf{66.36} & \textbf{63.38} & \textbf{85.76} & 93.60 & 4.23 & 1.218 \\
    sign\_64                    & 60.99 & 64.54 & 63.49 & 60.19 & 83.08 & 93.60 & 4.44 & 1.304 \\
    signed\_magnitude\_64       & 61.20 & 63.39 & 62.78 & 59.85 & 85.24 & \textbf{94.01} & 5.16 & \textbf{1.068} \\
    symmetry\_64                & 63.36 & 64.74 & 65.49 & 61.94 & 85.04 & 93.81 & 3.72 & 1.254 \\
    \midrule
    rcm\_binary\_64             & 59.03 & 62.82 & 62.04 & 58.30 & 81.01 & 90.71 & 5.99 & 1.304 \\
    rcm\_density\_64            & 61.51 & 65.59 & 64.19 & 61.00 & 84.31 & 93.91 & 4.23 & 1.289 \\
    rcm\_log\_density\_64       & 62.75 & 66.52 & 65.27 & 62.12 & 82.56 & 92.26 & 4.75 & 1.264 \\
    rcm\_magnitude\_64          & 63.47 & \textbf{66.80} & 65.74 & 62.79 & 85.35 & 93.81 & \textbf{3.41} & 2.766 \\
    rcm\_sign\_64               & 61.51 & 65.43 & 63.82 & 60.87 & 81.94 & 91.95 & 4.23 & 1.326 \\
    rcm\_signed\_magnitude\_64  & 62.54 & 65.17 & 64.09 & 61.36 & 83.90 & 92.67 & 4.44 & 1.069 \\
    \bottomrule
  \end{tabular}%
  }
\end{table}

Across the first two Campaigns where single image modes were used, the magnitude\_64 downsampling is the strongest single predictor.
It consistently outperforms its 13 competitors, winning Campaign~1 (Table~\ref{tab:results-campaign1}) with a 64.19\% accuracy and a 63.42\% F1 value and
Campaign~2 (Table~\ref{tab:results-campaign2}) with a 64.29\% accuracy and a 63.38\% F1 value.

While in both Campaigns the magnitude\_64 downsampling approach yields better results than their best baseline approaches, there are always some
configurations where nearly no improvement or even worse values can be observed.
In Campaign~1, the binary\_64 model, for example, has a lower accuracy and only an F1 value 0.03\% higher than the 1-NN baseline, showing that
the binary downsampling does not store that much valuable information when cropped to small images.
In Campaign~2, the binary\_64 model also does not show better results, but within this campaign the density\_64 approach was more notable as it was the only one
that dropped by over 8\%.
Investigation showed the model stopped predicting \texttt{minres+gamg} and \texttt{fcg+gamg} entirely, which resulted in a failure rate of 0\%.
The convergence penalty pushed the model away from those high-failure-rate classes.
By never predicting them it avoided failures, but at the cost of reducing their recall to 0\%.
This shows that a low failure rate alone is not a sufficient signal and must be read together with per-class recall, indicating that the convergence penalty as currently implemented is not sufficient. This will be investigated further in \autoref{sec:results-features}.

As only a limited number of downsampling methods made sense, to introduce additional variations, RCM reordering was applied to selected modes.
This had no significant impact in single-channel model results as values in most cases stayed within $\pm 0.6$\% of the original values.

The worst-performing modes consist of those that only regard structural features like binary and diagonal
and do not incorporate any magnitude information, confirming that entry values carry essential information for solver selection.
The six highest-ranking modes by F1 - \texttt{magnitude}, \texttt{rcm\_signed\_magnitude}, \texttt{rcm\_magnitude},
\texttt{signed\_magnitude}, \texttt{symmetry}, and \texttt{density} - are therefore carried forward as the basis
for the dual-channel experiments in Campaign~3~(\autoref{subsec:campaign3}).

The results also show that if a model with exceptional speed needs to be selected, in most cases it won't be the model with the
highest accuracy or the best F1 value.
In Campaign~1 for example, \texttt{magnitude\_64} achieves the best accuracy and F1, but \texttt{signed\_magnitude\_64} has a lower mean runtime ratio (1.069 vs.\ 1.259), making it the better choice when minimising runtime overhead is the priority.

\section{Impact of Feature Engineering and Convergence Penalty}\label{sec:results-features}

Table~\ref{tab:results-campaign2} (shown above in \autoref{sec:results-image-mode}) covers Campaign~2.
From Campaign~1 to Campaign~2, six additional matrix features and the convergence penalty~(\autoref{sec:training}) with $\lambda=1.0$ were introduced.
Because no isolated ablation was run for the 20-feature approach with $\lambda=0$ and the 14-feature approach with
$\lambda=1.0$, the effects of these changes need to be discussed simultaneously.
Note that from this point on, all experiments use the larger model~(\autoref{sec:architecture}).

\paragraph{Accuracy impact}
The accuracy gains from Campaign~1 to Campaign~2 are modest.
The nocnn MLP-only baseline approach improves by 1.44\% to 59.75\%, suggesting that the new features add useful
information.
Since the penalty primarily shapes which solver classes the model avoids rather than the feature representation itself, the
expanded feature set is the more likely driver of this particular gain.
For CNN-based models the improvement is negligible, the best approach being magnitude, only shifts 0.10\% to 64.29\% in accuracy.
This shows that information from the new features might already be present in some way within the CNN representations,
which would explain the diminishing results in the CNN model approaches vs.\ the nocnn approach.

\paragraph{Failure rate impact}
The metric where the convergence penalty is most clearly visible is the failure rate.
Combined failure rates across single-channel experiments dropped from 3.5\%--6.8\% in Campaign~1 (Table~\ref{tab:results-campaign1}) to
3.4\%--6.0\% in Campaign~2 (Table~\ref{tab:results-campaign2}).
This shows that the overall failure rate drops, sometimes significantly.
When looking at the per-class results for the magnitude downsampling with 14 features and without the penalty versus the
20 features and penalty approach, before the penalty, 4 classes had failure rates above 10\%, after, only 2 remain.
Unfortunately, for \texttt{symmlq+sor} the failure rate surged from 14.8\% to 28.2\%, which then drops
the overall failure rate back to previous levels.
A likely explanation is that the convergence penalty suppressed other high-failure classes, as discussed further below,
causing their probability mass to be redistributed towards solvers such as \texttt{symmlq+sor}.

This behaviour is reproduced in the other models as well, therefore the penalty approach needs to be examined as to
not only provide a shift in failure rates but realistic improvement.

\paragraph{Density side-effect}
The main negative side effect in this experiment is that the \texttt{density\_64} model collapsed in accuracy, dropping to 52.63\% from previously 60.27\%.
This loss of over 10\,\% is only observed in this model and might indicate a training instability, as the drop was not reproduced in later experiments,
but may also hint that the density encoding and the penalty gradient do not interact well.

\paragraph{Convergence penalty improvement}
As discussed in \autoref{sec:results-image-mode}, the convergence penalty can cause the model to stop predicting certain classes entirely,
reducing their recall to 0\% while also dropping their failure rate to 0\%.
This happens because the current penalty fires on every training sample and discourages the model from assigning any probability to solvers that
did not converge for that matrix.
For classes like \texttt{minres+gamg} that fail on many matrices, this accumulates through the training and the model stops predicting them entirely in the end.

Therefore, the convergence penalty needs improvement in future work.
Two promising directions that might help reduce the downsides of the current convergence penalty are the usage of
a threshold-gate, only penalising predictions when they exceed a certain probability, or incorporating the runtime ratio into
the penalty calculation.

\section{Single-Channel vs.\ Dual-Channel}\label{sec:results-single-vs-dual}

Adding a second CNN channel consistently improves over the single-channel approaches.
The best single-channel model, \texttt{magnitude\_64} (Table~\ref{tab:results-campaign2}: 64.29\% accuracy, 63.38\% F1), is outperformed by
\texttt{magnitude\_\_log\_density\_64} (Table~\ref{tab:results-campaign4}: 65.94\% accuracy, 65.02\% F1).
The best overall model is \texttt{magnitude\_\_rcm\_log\_density\_128} with 66.56\% accuracy, a 2.27\,\% improvement
over the best single-channel approach.

The size of the gain depends strongly on which two channels are chosen. When their information encodes different properties
they excel the most.
This can be seen by the combination of \texttt{magnitude} (absolute values) with \texttt{log\_density} (structural density) or \texttt{rcm\_log\_density} (density after bandwidth reduction).
Pairs with similar information like \texttt{rcm\_log\_density\_\_log\_density\_128} (two log-family channels) often drop even below some single-channel approaches
if their single-channel performance was already weak.

This confirms that a second channel is useful but only if the two downsampling methods contain non-redundant information.

Across all dual-channel experiments the most important channel is \texttt{magnitude}, as it consistently appears in the best
dual-channel pairs, reaffirming that magnitude encodes discriminative information essential for solver selection.


\subsection{Dual-Channel Results Overview}\label{subsec:dual-chan-results}

% ── Campaign 3: Dual-channel, 14 features, large model, 64px, no penalty ──────
\begin{table}[H]
  \centering
  \caption{Campaign~3 — Dual-channel: 14 features, no convergence penalty, large model, 64\,px.}
  \label{tab:results-campaign3}
  \resizebox{\textwidth}{!}{%
  \begin{tabular}{lrrrrrrrr}
    \toprule
    \textbf{Experiment} & \textbf{Acc\%} & \textbf{MP\%} & \textbf{MR\%} & \textbf{F1\%} & \textbf{Top-2\%} & \textbf{Top-3\%} & \textbf{Fail\%} & \textbf{MRT$\times$} \\
    \midrule
    1-NN                                           & 60.06 & 58.94 & 58.56 & 57.97 & 82.04 & 90.30 & 2.27 & 1.533 \\
    5-NN                                           & 62.95 & 62.46 & 61.85 & 61.12 & 82.56 & 90.40 & 3.20 & 1.605 \\
    nocnn                                          & 60.06 & 63.20 & 62.71 & 59.37 & 80.91 & 89.37 & 6.19 & 1.289 \\
    \midrule
    density\_\_rcm\_magnitude\_64                   & 64.40 & 67.28 & 65.93 & 63.72 & 85.66 & 93.40 & 5.06 & 1.219 \\
    density\_\_sign\_64                             & 63.98 & 66.90 & 66.42 & 63.14 & 85.35 & 94.12 & 4.44 & 1.217 \\
    density\_\_signed\_magnitude\_64               & 61.82 & 64.35 & 63.66 & 60.52 & 84.42 & 93.19 & 4.54 & 1.227 \\
    density\_\_symmetry\_64                        & 63.57 & 63.78 & 66.09 & 62.98 & 85.96 & 94.12 & 3.61 & 1.796 \\
    magnitude\_\_density\_64                        & 64.50 & 68.09 & 67.13 & 63.95 & \textbf{87.41} & 94.43 & 4.23 & 1.215 \\
    magnitude\_\_rcm\_magnitude\_64                 & 64.40 & 66.64 & 65.89 & 63.24 & 85.24 & 92.98 & 4.13 & 1.244 \\
    \textbf{magnitude\_\_rcm\_signed\_magnitude\_64}& \textbf{66.05} & \textbf{69.25} & 67.64 & 64.79 & 86.79 & 94.22 & 2.89 & 3.920 \\
    magnitude\_\_sign\_64                           & 64.50 & 66.79 & 66.93 & 63.52 & 85.96 & 93.91 & 4.23 & 1.215 \\
    \textbf{magnitude\_\_signed\_magnitude\_64}     & \textbf{66.05} & 68.21 & \textbf{68.23} & \textbf{65.32} & 86.48 & 94.12 & \textbf{2.79} & 4.476 \\
    magnitude\_\_symmetry\_64                       & 65.63 & 67.20 & 67.59 & 64.66 & 84.83 & 93.40 & 3.61 & 1.258 \\
    rcm\_magnitude\_\_sign\_64                      & 64.91 & 66.97 & 67.04 & 64.03 & 85.66 & 94.22 & 4.13 & 1.229 \\
    rcm\_magnitude\_\_signed\_magnitude\_64         & 65.63 & 68.90 & 67.61 & 64.86 & 85.45 & 94.22 & 3.82 & 1.235 \\
    rcm\_magnitude\_\_symmetry\_64                  & 64.29 & 66.85 & 66.74 & 63.34 & 85.55 & \textbf{94.53} & 3.61 & 3.298 \\
    rcm\_signed\_magnitude\_\_density\_64           & 63.36 & 65.94 & 65.10 & 61.89 & 84.83 & 92.67 & 4.54 & \textbf{1.063} \\
    rcm\_signed\_magnitude\_\_rcm\_magnitude\_64    & 63.88 & 66.95 & 65.97 & 63.22 & 85.66 & 93.81 & 4.02 & 1.221 \\
    rcm\_signed\_magnitude\_\_sign\_64              & 62.44 & 65.12 & 64.14 & 60.79 & 84.93 & 93.70 & 4.33 & 1.226 \\
    rcm\_signed\_magnitude\_\_signed\_magnitude\_64 & 62.33 & 64.34 & 63.62 & 60.55 & 84.11 & 93.50 & 4.64 & 1.065 \\
    rcm\_signed\_magnitude\_\_symmetry\_64          & 62.85 & 63.05 & 64.26 & 61.81 & 83.49 & 92.26 & 3.82 & 1.631 \\
    signed\_magnitude\_\_sign\_64                   & 62.33 & 64.70 & 63.94 & 61.29 & 84.00 & 93.60 & 4.64 & 1.235 \\
    symmetry\_\_sign\_64                            & 63.05 & 65.83 & 65.22 & 61.92 & 85.66 & 93.91 & 3.72 & 3.093 \\
    symmetry\_\_signed\_magnitude\_64               & 63.36 & 64.72 & 64.92 & 62.08 & 85.96 & \textbf{94.53} & 3.20 & 1.930 \\
    \bottomrule
  \end{tabular}%
  }
\end{table}

% ── Campaign 4: Dual-channel, 20 features, convergence penalty λ=1.0 ──────────
\begin{table}[H]
  \centering
  \caption{Campaign~4 — Dual-channel: 20 features, convergence penalty $\lambda=1.0$, large model, 64\,px and 128\,px.}
  \label{tab:results-campaign4}
  \resizebox{\textwidth}{!}{%
  \begin{tabular}{lrrrrrrrr}
    \toprule
    \textbf{Experiment} & \textbf{Acc\%} & \textbf{MP\%} & \textbf{MR\%} & \textbf{F1\%} & \textbf{Top-2\%} & \textbf{Top-3\%} & \textbf{Fail\%} & \textbf{MRT$\times$} \\
    \midrule
    1-NN                                              & 59.44 & 59.02 & 58.34 & 57.00 & 82.15 & 90.40 & 1.86 & 1.539 \\
    5-NN                                              & 63.57 & 63.25 & 62.38 & 61.61 & 82.97 & 91.02 & 2.89 & 1.589 \\
    nocnn                                             & 59.75 & 62.40 & 62.72 & 59.08 & 80.80 & 89.78 & 4.95 & 2.926 \\
    \midrule
    magnitude\_\_log\_density\_64                     & 65.94 & 68.42 & 67.73 & 65.02 & 86.79 & 94.63 & 4.02 & 1.214 \\
    magnitude\_\_log\_density\_128                    & 65.94 & 68.67 & 68.22 & \textbf{65.60} & 87.00 & 94.63 & 4.13 & 1.051 \\
    magnitude\_\_rcm\_log\_density\_64                & 65.43 & 67.72 & 67.75 & 64.74 & 86.27 & 94.63 & 4.02 & 1.220 \\
    \textbf{magnitude\_\_rcm\_log\_density\_128}      & \textbf{66.56} & 69.15 & \textbf{68.88} & 65.53 & 86.38 & 94.53 & 3.51 & 1.051 \\
    magnitude\_\_rcm\_magnitude\_64                   & 63.78 & 66.49 & 65.76 & 63.37 & 86.27 & 94.63 & 3.92 & 1.227 \\
    magnitude\_\_rcm\_magnitude\_128                  & 65.02 & 66.73 & 66.62 & 64.22 & 87.10 & 94.74 & 3.61 & 1.222 \\
    magnitude\_\_rcm\_signed\_magnitude\_64           & 65.53 & 67.52 & 67.36 & 64.30 & 85.96 & 93.40 & 3.72 & 1.234 \\
    magnitude\_\_rcm\_signed\_magnitude\_128          & 65.84 & \textbf{69.39} & 67.59 & 65.16 & 86.38 & 94.43 & 3.72 & 1.213 \\
    magnitude\_\_symmetry\_64                         & 64.71 & 66.34 & 66.35 & 63.49 & 85.45 & 94.01 & \textbf{2.68} & 3.821 \\
    magnitude\_\_symmetry\_128                        & 64.60 & 66.86 & 66.41 & 63.87 & 85.55 & 93.81 & 3.61 & 1.086 \\
    rcm\_log\_density\_\_log\_density\_64             & 62.02 & 66.30 & 65.07 & 61.44 & 82.66 & 92.47 & 4.23 & 1.276 \\
    rcm\_log\_density\_\_log\_density\_128            & 63.26 & 66.32 & 65.71 & 62.28 & 84.73 & 94.32 & 4.64 & 1.223 \\
    rcm\_magnitude\_\_log\_density\_64                & 64.09 & 67.24 & 66.56 & 63.77 & 85.66 & 93.29 & 3.92 & 1.247 \\
    rcm\_magnitude\_\_log\_density\_128               & 64.60 & 67.97 & 67.00 & 64.35 & 86.17 & 93.91 & 4.02 & 1.227 \\
    rcm\_magnitude\_\_rcm\_log\_density\_64           & 63.98 & 66.95 & 65.98 & 63.51 & 85.35 & 93.40 & 3.82 & 1.256 \\
    rcm\_magnitude\_\_rcm\_log\_density\_128          & 64.50 & 68.06 & 67.03 & 64.19 & 84.52 & 93.19 & 3.51 & 1.282 \\
    rcm\_magnitude\_\_rcm\_signed\_magnitude\_64      & 64.29 & 67.98 & 66.25 & 63.50 & 85.55 & 94.53 & 4.02 & 1.218 \\
    rcm\_magnitude\_\_rcm\_signed\_magnitude\_128     & 64.60 & 67.94 & 66.77 & 64.30 & 86.79 & 94.32 & 4.02 & \textbf{1.048} \\
    rcm\_magnitude\_\_symmetry\_64                    & 63.67 & 64.68 & 65.37 & 62.66 & 86.17 & 94.43 & 3.20 & 2.765 \\
    rcm\_magnitude\_\_symmetry\_128                   & 65.02 & 66.05 & 66.83 & 64.23 & \textbf{87.20} & 94.84 & 2.99 & 3.063 \\
    rcm\_signed\_magnitude\_\_log\_density\_64        & 63.47 & 66.25 & 65.69 & 62.40 & 85.66 & 93.60 & 4.23 & 1.057 \\
    rcm\_signed\_magnitude\_\_log\_density\_128       & 63.36 & 66.43 & 65.51 & 62.54 & 86.27 & 94.53 & 4.13 & 1.228 \\
    rcm\_signed\_magnitude\_\_rcm\_log\_density\_64   & 64.40 & 66.05 & 65.84 & 63.42 & 85.96 & 94.12 & 3.92 & 1.072 \\
    rcm\_signed\_magnitude\_\_rcm\_log\_density\_128  & 63.67 & 66.50 & 65.11 & 62.52 & 85.86 & 94.01 & 4.02 & 1.263 \\
    symmetry\_\_log\_density\_64                      & 64.50 & 66.70 & 66.88 & 63.57 & 87.10 & 95.15 & 3.10 & 2.540 \\
    symmetry\_\_log\_density\_128                     & 65.12 & 66.68 & 67.14 & 64.10 & 86.79 & 94.84 & 3.10 & 3.336 \\
    symmetry\_\_rcm\_log\_density\_64                 & 64.29 & 67.12 & 66.69 & 63.43 & 85.66 & 94.01 & 3.10 & 4.058 \\
    symmetry\_\_rcm\_log\_density\_128                & 64.71 & 66.92 & 66.98 & 63.84 & 86.58 & \textbf{95.36} & 2.99 & 2.126 \\
    symmetry\_\_rcm\_signed\_magnitude\_64            & 64.09 & 65.96 & 65.58 & 62.62 & 84.31 & 93.29 & 3.20 & 1.817 \\
    symmetry\_\_rcm\_signed\_magnitude\_128           & 64.09 & 65.31 & 66.05 & 62.89 & 86.69 & 95.15 & 3.10 & 2.710 \\
    \bottomrule
  \end{tabular}%
  }
\end{table}

The dual-channel experiments are summarised in Table~\ref{tab:results-campaign3} (Campaign~3, 14 features, 64\,px)
and Table~\ref{tab:results-campaign4} (Campaign~4, 20 features, 64\,px and 128\,px).
For Campaign~3 the six best image modes from Campaign~1 - \texttt{magnitude}, \texttt{rcm\_signed\_magnitude},
\texttt{rcm\_magnitude}, \texttt{signed\_magnitude}, \texttt{symmetry}, and \texttt{density} -
together with \texttt{sign} are combined into all 21 pairwise combinations.
For Campaign~4 only the six best image modes from Campaign~2 - \texttt{magnitude}, \texttt{rcm\_magnitude},
\texttt{rcm\_log\_density}, \texttt{symmetry}, \texttt{log\_density}, and \texttt{rcm\_signed\_magnitude} - are combined,
yielding 15 pairs each evaluated at two resolutions (64\,px and 128\,px).

The dominant finding over these two Campaigns is that \texttt{magnitude} is the most informative channel in the dual-channel setting as well.
Every top-performing pair across both campaigns includes it as one channel.
Replacing \texttt{magnitude} with any other mode as the primary channel reduces accuracy by 2--3\,\%.

The best second channel is \texttt{log\_density} (best F1) or \texttt{rcm\_log\_density}
(best accuracy), both outperforming \texttt{symmetry} and \texttt{rcm\_signed\_magnitude}.
The combination \texttt{rcm\_log\_density+log\_density} (two log-family channels) performs
worst in Campaign~4, confirming that redundant modes do not benefit the model.

128\,px consistently outperforms 64\,px by $+0.5$--$+1.3$\,\% across all mode pairs,
at a cost of roughly $4\times$ longer training time and $4\times$ higher inference latency.
The best Campaign~4 single model is \texttt{magnitude+rcm\_log\_density} at 128\,px
with 66.56\% accuracy and 65.53\% F1.

Dual-channel models outperform the best single-channel model by $+2.3$\,\% in accuracy
(66.56\% vs.\ 64.29\%), confirming that the two image channels carry complementary
discriminative information beyond what either channel provides alone.

\section{Effect of Image Resolution}\label{sec:results-resolution}

In the dual-channel experiment with 64 and 128\,px images, shown in \autoref{tab:results-campaign4}, the increase in resolution
brings a consistent improvement of 0.5--1.3\,\% in accuracy and F1 values across all 15 mode pairs.
The largest gain is seen in \texttt{magnitude+rcm\_log\_density}, where accuracy improves by 1.13\,\% and F1 by 0.79\,\%.
The smallest gain is on \texttt{magnitude+log\_density}, where accuracy stays the same and only F1 increases by 0.58\,\%.
Almost all pairs improve or remain stable at higher resolution.

The cost of the higher resolution is a $4\times$ increase in inference time, as the dual-channel 64\,px model runs at approximately 2.8\,ms/sample and the
128\,px model takes about 10.9\,ms/sample.
The training time per model is around 13\,min at 64\,px and around 48\,min at 128\,px, on a NVIDIA GeForce RTX 5060 Laptop GPU.

In the single-channel experiments (Table~\ref{tab:results-campaign1}), results were more mixed.
\texttt{binary\_128} and \texttt{density\_128} outperformed their 64\,px counterparts, but for the magnitude mode the 64\,px version outperformed the 128\,px version in accuracy and F1.
This suggests that 64\,px is already sufficient for magnitude-encoded images, where per-cell absolute values matter more than the
fine spatial arrangement of non-zeros.

\section{Ensemble vs.\ Single Best Model}\label{sec:results-ensemble}

% ── Ensembles: Campaign 3 and Campaign 4 ─────────────────────────────────────
\begin{table}[H]
  \centering
  \caption{Campaign~5 - Ensemble results: averaging of four independently trained dual-channel models.
           \emph{Campaign~3}: 14 features, 64\,px, no penalty. Members are
           \texttt{mag+signed\_mag}, \texttt{mag+rcm\_signed\_mag},
           \texttt{mag+symmetry}, \texttt{rcm\_mag+signed\_mag}.
           \emph{Campaign~4}: 20 features, 128\,px, $\lambda=1.0$. Members are
           \texttt{mag+log\_density}, \texttt{mag+rcm\_log\_density},
           \texttt{rcm\_mag+rcm\_signed\_mag}, \texttt{symmetry+log\_density}.
           Bold marks the column best.}
  \label{tab:results-campaign5}
  \resizebox{\textwidth}{!}{%
  \begin{tabular}{lrrrrrrrr}
    \toprule
    \textbf{Ensemble} & \textbf{Acc\%} & \textbf{MP\%} & \textbf{MR\%} & \textbf{F1\%} & \textbf{Top-2\%} & \textbf{Top-3\%} & \textbf{Fail\%} & \textbf{MRT$\times$} \\
    \midrule
    Campaign~3 ensemble (64\,px, 4 models) & \textbf{67.39} & 69.47 & \textbf{69.58} & 66.37 & \textbf{87.82} & \textbf{94.94} & \textbf{2.79} & 4.055 \\
    Campaign~4 ensemble (128\,px, 4 models) & 67.29 & \textbf{70.37} & 69.48 & \textbf{66.77} & 87.41 & \textbf{94.94} & 3.92 & \textbf{1.043} \\
    \bottomrule
  \end{tabular}%
  }
\end{table}

\paragraph[Campaign 3 ensemble.]{Campaign~3 ensemble.}

The Campaign~3 ensemble (all 64\,px):
\begin{itemize}
  \item \texttt{magnitude+signed\_magnitude}
  \item \texttt{magnitude+rcm\_signed\_magnitude}
  \item \texttt{magnitude+symmetry}
  \item \texttt{rcm\_magnitude+signed\_magnitude}
\end{itemize}
achieves 67.39\% accuracy and an F1 of 66.37\%, surpassing the best individual member
(\texttt{magnitude+signed\_magnitude\_64} at 66.05\% accuracy and 65.32\% F1)
by $+1.34$\,\% in accuracy and $+1.05$\,\% in F1.
The failure rate drops to 2.79\%, below all four member models individually,
confirming that non-converging predictions are not correlated across members.

\paragraph[Campaign 4 ensemble.]{Campaign~4 ensemble.}
The Campaign~4 ensemble (all 128\,px):
\begin{itemize}
  \item \texttt{magnitude+log\_density}
  \item \texttt{magnitude+rcm\_log\_density}
  \item \texttt{rcm\_magnitude+rcm\_signed\_magnitude}
  \item \texttt{symmetry+log\_density}
\end{itemize}
achieves 67.29\% accuracy and 66.77\% F1, compared to the best individual member
\texttt{magnitude+rcm\_log\_density\_128} (66.56\% accuracy and 65.53\% F1),
a gain of $+0.73$\,\% in accuracy and $+1.24$\,\% in F1.
The failure rate of 3.92\% and mean runtime ratio of 1.043$\times$ remain below or comparable to all individual members.

Notably, the Campaign~4 ensemble surpasses the Campaign~3 ensemble only in F1 (66.77\% vs.\ 66.37\%),
while trailing in accuracy (67.29\% vs.\ 67.39\%).
This comes at notably higher inference cost, running at approximately 46.5\,ms/sample —
roughly $4\times$ slower than the Campaign~3 ensemble ($\approx$12\,ms) — due to the 128\,px images.

\paragraph{Inference cost.}
The ensemble multiplies inference time by the number of members ($4\times$).
At $\approx$12\,ms/sample for four dual-channel 64\,px models (Campaign~3) and
$\approx$46.5\,ms/sample for four dual-channel 128\,px models (Campaign~4), the
prediction overhead remains negligible compared to typical iterative solver runtimes
which range from tens of milliseconds to several seconds per solve.

\section{Comparison to MM-AutoSolver}\label{sec:comparison-mm}

This section compares the thesis model against the MM-AutoSolver paper results and the replication performed as a baseline.
The replication was trained on a synthetic dataset with some real SuiteSparse matrices included, but was
evaluated on the same 19 solver classes with slightly different per-class distributions.
The replication even exceeds the paper in accuracy by 3.43\,\%, as seen in \autoref{tab:comparison-mm}, and
may be the result of the slightly different data source~(\autoref{sec:dataset}).
\begin{table}[H]
  \centering
  \caption{Comparison against MM-AutoSolver~\cite{xiong2025mmautosolver}.
           Best model: \texttt{magnitude+rcm\_log\_density} at 128\,px (Campaign~4).
           Ensemble C3: 4 models, 64\,px, 14 features.
           Ensemble C4: 4 models, 128\,px, 20 features, $\lambda=1.0$.
           Bold marks the best thesis result in each column.}
  \label{tab:comparison-mm}
  \begin{tabular}{lrrrr}
    \toprule
    \textbf{System} & \textbf{Acc\%} & \textbf{MP\%} & \textbf{MR\%} & \textbf{F1\%} \\
    \midrule
    MM-AutoSolver (paper)              & 78.54 & 63.41 & 62.81 & 62.53 \\
    MM-AutoSolver (replication)        & 81.97 & 62.88 & 61.83 & 60.45 \\
    \midrule
    This thesis (best model, C4)      & 66.56 & 69.15 & 68.88 & 65.53 \\
    This thesis (ensemble C3, 64\,px)  & \textbf{67.39} & 69.47 & \textbf{69.58} & 66.37 \\
    This thesis (ensemble C4, 128\,px) & 67.29 & \textbf{70.37} & 69.48 & \textbf{66.77} \\
    \bottomrule
  \end{tabular}
\end{table}

A direct comparison between the MM-AutoSolver approach and the thesis approach is misleading for two reasons.
First, the class distribution in the dataset differs substantially from the MM-AutoSolver dataset, where the \texttt{fbcgsr+jacobi} and \texttt{bcgsl+none} classes make up more than 36\% of the whole distribution.
A model can exploit this by defaulting to one of these dominant classes when uncertain, inflating accuracy.
The thesis dataset was constructed with a minimal per-class constraint of 250 matrices and capped at 600 matrices per class,
resulting in a much flatter distribution and harder problem.
Second, the training matrices also differ, as the MM-AutoSolver paper uses more real-world application matrices generated
within FreeFEM~\cite{hecht2012freefem} and OpenFOAM~\cite{weller1998tensorial}.
In this thesis, only SuiteSparse matrices with real-world applications are used, while the remainder consists of synthetically generated ones.

Despite the lower accuracy, the macro-averaged metrics are less sensitive to class distribution and can still be compared.
The thesis model even achieves around 7\,\% better results in Macro Precision and Macro Recall, meaning that despite the harder
problem it learned to distinguish the classes better than the MM-AutoSolver approach.

Another insight is the Top-3 accuracy of around 94\%, meaning that even when the predicted solver is not the best,
in nearly all cases it is among the top three choices.
Combined with an average failure rate of around 3--4\%, the approach is practically viable.

\section{Per-Solver Analysis}\label{sec:per-solver}

\autoref{tab:per-solver-ensemble} breaks down per-class F1, precision, recall, and failure rate
for the Campaign~4 ensemble (members: \texttt{magnitude+log\_density},
\texttt{magnitude+rcm\_log\_density}, \texttt{rcm\_magnitude+rcm\_signed\_magnitude},
\texttt{symmetry+log\_density}; 20 features, 128\,px, $\lambda=1.0$; 67.29\% Acc, 66.77\% F1).
The 19 classes span a wide performance range, from \texttt{minres+gamg} at 88.7\% F1 down to
\texttt{gmres+gamg} at 20.9\% F1, illustrating that solver prediction difficulty varies
considerably across the label space.


\begin{table}[H]
  \centering
  \caption[Per-class results for the Campaign~4 ensemble.]{Per-class results for the Campaign~4 ensemble (4 models: magnitude+log\_density,
           magnitude+rcm\_log\_density, rcm\_magnitude+rcm\_signed\_magnitude,
           symmetry+log\_density, 128\,px, $\lambda=1.0$).
           Sorted by F1 descending. N is the number of training samples per class.}
  \label{tab:per-solver-ensemble}
  \begin{tabular}{lrrrrr}
    \toprule
    \textbf{Solver pair} & \textbf{N} & \textbf{F1\%} & \textbf{Prec\%} & \textbf{Rec\%} & \textbf{Fail\%} \\
    \midrule
    minres+gamg      &  600 & 88.7 & 80.8 & 98.3 &  0.0 \\
    cr+eisenstat     &  600 & 88.3 & 96.1 & 81.7 &  0.0 \\
    symmlq+sor       &  250 & 86.3 & 84.6 & 88.0 & 15.4 \\
    symmlq+jacobi    &  253 & 82.1 & 74.2 & 92.0 &  3.2 \\
    bcgsl+asm        &  428 & 73.4 & 78.4 & 69.0 &  2.7 \\
    fbcgsr+jacobi    &  600 & 72.3 & 72.9 & 71.7 &  1.7 \\
    cr+ilu           &  600 & 72.1 & 71.0 & 73.3 &  6.5 \\
    cr+jacobi        &  600 & 71.9 & 75.9 & 68.3 &  0.0 \\
    cg+bjacobi       &  600 & 71.9 & 67.6 & 76.7 &  4.4 \\
    cg+eisenstat     &  600 & 70.7 & 73.2 & 68.3 &  0.0 \\
    bcgsl+none       &  600 & 70.0 & 70.0 & 70.0 &  0.0 \\
    dgmres+none      &  354 & 66.7 & 64.9 & 68.6 &  0.0 \\
    symmlq+icc       &  600 & 63.8 & 56.4 & 73.3 &  5.1 \\
    fbcgsr+ilu       &  600 & 63.4 & 78.0 & 53.3 &  0.0 \\
    cgs+gamg         &  336 & 52.8 & 38.4 & 84.8 &  0.0 \\
    fcg+gamg         &  290 & 52.3 & 39.0 & 79.3 &  3.4 \\
    cg+ilu           &  600 & 52.1 & 52.5 & 51.7 & 30.5 \\
    fgmres+gamg      &  600 & 49.0 & 63.2 & 40.0 &  0.0 \\
    gmres+gamg       &  600 & 20.9 &100.0 & 11.7 &  0.0 \\
    \bottomrule
  \end{tabular}
\end{table}

\paragraph{Easy classes.}
\texttt{minres+gamg} (88.7\% F1, recall 98.3\%) leads the ranking, reliably identifying
symmetric indefinite problems that are unsuitable for CG-type solvers.
\texttt{cr+eisenstat} (88.3\% F1, precision 96.1\%) follows closely, dominating structurally
symmetric and nearly diagonally dominant matrices.
\texttt{symmlq+jacobi} (82.1\% F1, recall 92.0\%) shows high recall but lower precision
(74.2\%), indicating the model confidently identifies this class but occasionally predicts
it for matrices belonging to neighbouring symmetric solver classes.

\paragraph{Hard classes.}
\texttt{gmres+gamg} is the hardest class (20.9\% F1) with a recall of only 11.7\%
despite perfect precision of 100.0\%: the model is always right when it predicts this class,
but almost never does so.
This is a consequence of class ambiguity within the GAMG-preconditioned family:
\texttt{fgmres+gamg} (49.0\% F1) and \texttt{fcg+gamg} (52.3\% F1) converge in similar
time on many of the same matrices, and the model spreads probability across the family
rather than committing to one member.
\texttt{fbcgsr+ilu} (63.4\% F1, recall 53.3\%) and \texttt{symmlq+icc} (63.8\% F1) are
the next weakest: both are specialised combinations whose winning conditions are narrow
and hard to separate from neighbouring classes by image alone.

\paragraph{Failure rates.}
\texttt{cg+ilu} stands out with a 30.5\% failure rate, while \texttt{symmlq+sor} (15.4\%) and \texttt{cr+ilu} (6.5\%) also carry elevated rates.
All three rely on preconditioners sensitive to matrix conditioning, making them harder to predict reliably.
Despite these per-class spikes, the overall ensemble failure rate remains low at 3.92\%.

\section{Practical Inference Cost}\label{sec:inference-cost}

Model inference is fast enough to be negligible compared to actual solver runtimes.
\autoref{tab:inference-latency} summarises the per-sample CPU latency for the main model configurations, measured on the
969 validation samples.

\begin{table}[h]
  \centering
  \caption{CPU inference latency per sample (averaged over 969 validation samples).}
  \label{tab:inference-latency}
  \begin{tabular}{lrr}
    \toprule
    \textbf{Configuration} & \textbf{Latency (ms/sample)} & \textbf{Model params} \\
    \midrule
    MLP only (nocnn)                              & $\approx$0.08 & $\approx$57\,K \\
    Single-channel CNN, small model (64\,px)      & $\approx$1.2  & $\approx$0.5\,M \\
    Single-channel CNN, small model (128\,px)     & $\approx$4.5  & $\approx$0.5\,M \\
    Single-channel CNN, large model (64\,px)      & $\approx$2.8  & $\approx$2.1\,M \\
    Dual-channel CNN, large model (64\,px)        & $\approx$3.0  & $\approx$2.1\,M \\
    Dual-channel CNN, large model (128\,px)       & $\approx$10.9 & $\approx$2.1\,M \\
    Campaign~3 ensemble (64\,px, 4 models)        & $\approx$12.0 & $4\times$2.1\,M \\
    Campaign~4 ensemble (128\,px, 4 models)       & $\approx$46.5 & $4\times$2.1\,M \\
    \bottomrule
  \end{tabular}
\end{table}

The MLP-only baseline runs at under 0.1\,ms per sample, confirming that the scalar feature branch alone imposes negligible overhead.
Adding the CNN increases latency to around 1--3\,ms for 64\,px models, while 128\,px models are roughly $4\times$ slower, consistent with the fourfold increase in pixel count.
The dual-channel configuration adds only marginal cost over the equivalent single-channel model (2.8\,ms vs.\ 3.0\,ms at 64\,px), meaning the second image channel comes at almost no additional inference expense.

\paragraph{Full pipeline cost.}
Model inference alone does not capture the full overhead: feature extraction and image rendering must also be accounted for.
To illustrate the combined cost, a small timing experiment is conducted on two representative matrix types, a 2D Poisson grid (SPD) and a random non-symmetric matrix, each at two sizes, $n \approx 50\text{k}$ and $n \approx 100\text{k}$.
This is not a systematic benchmark but a targeted measurement to show where the pipeline time is spent.
The non-symmetric matrices use 50 non-zeros per row, giving $\text{nnz} \approx 2.6\text{M}$ and $\text{nnz} \approx 5.1\text{M}$ respectively, compared to $\approx 250\text{k}$ and $\approx 500\text{k}$ for the Poisson grids.
All 19 solvers converge within 10\,s on the SPD matrices and 11 of 19 converge within 10\,s on each non-symmetric matrix, with the remaining 8 either failing to converge or running beyond the 10\,s kill threshold set for this experiment.
The random baseline simulates picking uniformly at random until a solver converges: each failed attempt costs its detection time (capped at 10\,s), and the expected number of failures before success scales with $n_\text{failed}/n_\text{success}$.

Table~\ref{tab:pipeline-breakdown} shows the complete breakdown for the best-performing models from Campaigns~3 and~4.

\begin{table}[H]
  \centering
  \caption{
    Full pipeline breakdown for \texttt{rcm\_signed\_magnitude\_\_density\_64} (C3, MRT$\times$\,=\,1.063)
    and \texttt{magnitude\_\_rcm\_log\_density\_128} (C4, MRT$\times$\,=\,1.051) across four test matrices.
    All times in ms; image rendering covers both channels.
    Solvers are killed after 10\,s; the random baseline uses a geometric retry model.%
  }
  \label{tab:pipeline-breakdown}
  \resizebox{\textwidth}{!}{
  \begin{tabular}{lrrrrrrrr}
    \toprule
    & \multicolumn{4}{c}{\textbf{C3: rcm\_signed\_magnitude\_\_density\_64 (14 feat, 64\,px)}}
    & \multicolumn{4}{c}{\textbf{C4: magnitude\_\_rcm\_log\_density\_128 (20 feat, 128\,px)}} \\
    \cmidrule(lr){2-5}\cmidrule(lr){6-9}
    & \textbf{SPD 50k} & \textbf{SPD 100k} & \textbf{NS 50k} & \textbf{NS 100k}
    & \textbf{SPD 50k} & \textbf{SPD 100k} & \textbf{NS 50k} & \textbf{NS 100k} \\
    \midrule
    Feature extraction (ms)      &   33 &   23 &    93 & 384    &   521 &   600 &  2{,}617 & 5{,}424 \\
    Image rendering, 2ch.\ (ms)  &   29 &   55 &   304 & 641    &    28 &    54 &    302 &   631 \\
    Model inference (ms)         &    3 &    3 &     3 &   3    &    11 &    11 &     11 &    11 \\
    \midrule
    \textbf{Pipeline overhead (ms)} & \textbf{65} & \textbf{81} & \textbf{400} & \textbf{1{,}028} & \textbf{560} & \textbf{665} & \textbf{2{,}930} & \textbf{6{,}066} \\
    Predicted solve, MRT$\times$ (ms) & 108 & 237 &  39 &  55   &  107 &  234 &     38 &    54 \\
    \textbf{Total (ms)}          & \textbf{173} & \textbf{318} & \textbf{439} & \textbf{1{,}083} & \textbf{667} & \textbf{899} & \textbf{2{,}968} & \textbf{6{,}120} \\
    \midrule
    Best solver (ms)             & 102 & 223 &  37 &  51   &  102 &  223 &     37 &    51 \\
    Random baseline (ms)         & 270 & 1{,}225 & 6{,}468 & 7{,}484 & 270 & 1{,}225 & 6{,}468 & 7{,}484 \\
    \textbf{Speedup vs.\ random} & \textbf{1.56$\times$} & \textbf{3.85$\times$} & \textbf{14.75$\times$} & \textbf{6.91$\times$}
                                 & \textbf{0.40$\times$} & \textbf{1.36$\times$} & \textbf{2.18$\times$} & \textbf{1.22$\times$} \\
    \bottomrule
  \end{tabular}%
  }
\end{table}

The dominant cost in the 20-feature set is feature~16, the structural symmetry fraction, which checks for matching $(i,j)$--$(j,i)$ non-zero pairs using set operations that scale with $\text{nnz}$.
It accounts for 506\,ms on the sparse SPD 50k matrix, 589\,ms on SPD 100k, and 2.5\,s and 5.2\,s on the two non-symmetric matrices, making it responsible for the vast majority of C4's extraction time in every case.
The 14-feature set avoids this computation entirely, giving C3 a decisive pipeline advantage.

C3 achieves a positive speedup over random selection in all four cases, ranging from 1.56$\times$ on the well-conditioned SPD 50k matrix to 14.75$\times$ on the non-symmetric 50k matrix where 8 of 19 solvers fail or are killed.
C4 is slower than random on the small SPD matrix (0.40$\times$), where the 561\,ms overhead alone exceeds the 270\,ms random baseline, but achieves a 1.36$\times$ speedup on the larger one, where solver runtimes are long enough to offset the 664\,ms overhead.
On non-symmetric matrices C4 recovers a modest speedup (2.18$\times$ and 1.22$\times$), but remains 5--7$\times$ slower than C3 in absolute terms.

Both models were selected for their low MRT$\times$ values (C3: 1.063, C4: 1.051), meaning the predicted solver is only marginally slower than optimal on average, so solver-selection quality is comparable between the two campaigns.
These results show that the 14-feature pipeline of Campaign~3 offers a better cost-quality trade-off across all tested matrix types and sizes as it consistently beats random selection and keeps the pipeline overhead below 1.1\,s even at $n = 100\text{k}$.
Campaign~4's 20-feature pipeline is only competitive when expected solver runtimes substantially exceed its extraction overhead, which itself scales with matrix density and size.
